\documentclass[../../main.tex]{subfiles}
\begin{document}

\begin{flushright}
\footnotesize\textit{【自動文字起こし・要確認】}
\end{flushright}

[解] $\alpha$と$l$の交点を$H$とする。$H$は$z=0$上で
$y = -x$ と $y = x - 1$ の交点で、$H(\dfrac{1}{2}, -\dfrac{1}{2}, 0)$ となる。

したがって、$t \ge \dfrac{1}{2}$ の時、$P$と$B$、$P$と$C$は互いに$\alpha$に
関して反対側にあり、$PB, PC$は$\alpha$と交点$D, E$を持つ。ここで、$\alpha$を$y$軸だから、右下図で
相似から、
\[ \overline{EH} = \dfrac{t - \dfrac{1}{2}}{t} \cdots \text{①} \]
である。

次に$D$にも同様に
\[ \overline{AD} = \dfrac{1}{1 + (t - \dfrac{1}{2})} \overline{AH} \]
\[ = \dfrac{1}{t + \dfrac{1}{2}} \dfrac{\sqrt{2}}{2} \cdots \text{②} \]

したがって、
\[ S(t) = (\triangle ADE \text{の面積}) \]
\[ = \dfrac{1}{2} \overline{AD} \cdot \overline{EH} \]
\[ = \dfrac{1}{2} \dfrac{1}{t + \dfrac{1}{2}} \dfrac{\sqrt{2}}{2} \dfrac{t - \dfrac{1}{2}}{t} \quad (\because \text{①②}) \]
\[ = \dfrac{\sqrt{2}}{4} \dfrac{p}{(p + 1)(p + \dfrac{1}{2})} \quad (p = t - \dfrac{1}{2} \ge 0) \]
\[ = \dfrac{\sqrt{2}}{4} \dfrac{1}{p + \dfrac{1}{2p} + \dfrac{3}{2}} \le \dfrac{\sqrt{2}}{4} \dfrac{1}{2\sqrt{\dfrac{1}{2}} + \dfrac{3}{2}} \quad (\because \text{AM-GM}, p > 0) \]
\[ = \dfrac{1}{2}(-4 + 3\sqrt{2}) \]
等号成立は $p = \dfrac{1}{2p} \therefore p = \dfrac{\sqrt{2}}{2} \ (>0)$の時である。\hfill(終)

\begin{center}
\begin{tikzpicture}
% A generic representation of the figures in the handwritten document
\draw[->] (0,0) -- (2,0) node[right] {$y$};
\draw[->] (0,0) -- (-1,-1) node[below left] {$x$};
\draw[->] (0,0) -- (0,2) node[above] {$z$};
\draw[dashed] (-0.5,-0.5) -- (0.5,1.5) -- (1.5,0.5) -- cycle;
\node at (0.5,1.5) [above right] {$C$};
\node at (1.5,0.5) [below right] {$B$};
\node at (-0.5,-0.5) [below left] {$A$};
\end{tikzpicture}
\end{center}

\end{document}
